\section{Matematisk \& teknisk grund}
\cvitem{}{
Från civilingenjörsutbildningen i datateknik (Linköpings universitet, 180 hp, inriktning mjukvaruteknik)
kommer följande kurser som direkt ligger till grund för undervisning i matematik och fysik:
}
\cvitem{Matematik (66 hp)}{
Analys (en och flera variabler), linjär algebra, diskret matematik och logik, sannolikhetslära, statistik,
numeriska algoritmer, kombinatorisk optimering, transformteori, signalteori
}
\cvitem{Elektronik \& fysik (55 hp)}{
Grundläggande elektronik och mätteknik, kretsteori, switchteori och logisk design, datorarkitektur och
hårdvara, digitalt projektlabb, mikrodatorprojektlabb, elektromagnetism, modern fysik, reglerteknik
}
\cvitem{Programmering \& tillämpad logik (45 hp)}{
Datastrukturer och algoritmer, objektorienterad utveckling, samtidig programmering och operativsystem,
datorgrafik, artificiell intelligens, databasteknik, datasäkerhet
}
\cvitem{Examensarbete}{
\emph{Evolving an OS Kernel Using Temporal Logic and Aspect-Oriented Programming} -- tillämpning av
formell (temporal) logik för verifiering av mjukvara, vid École des Mines de Nantes, Frankrike.
}
\cvitem{}{
Tidigare, på gymnasienivå (Naturvetenskapligt program, Rönneskolan), ingick matematik, fysik och kemi,
tillsammans med ett tekniskt tillval i programmering och elektronik. Ett specialarbete byggde 3D-/vektorgrafik
i assembler på Amigan, med tillämpning av rotationsmatriser, linjeritningsalgoritmer och
koordinattransformationer -- linjär algebra omsatt direkt i visuell, handgriplig praktik.
}
